\Needspace{75pt}
\pdfbookmark[1]{Calculation of the average values of contributions and pensions}{witwen-mittelwerthe}
\begin{center}[\textit{Calculation of the average values of contributions and pensions.}\end{center}
\noindent\hspace*{\parindent}The number of members who, in an uninterrupted sequence according to the numbering of their entry dates, had died or otherwise left the professors' widows' fund by 1 October 1851 is 134. For each of these 134 former members, the following table gives the year in which entry took place on 1 Oct. (and only exceptionally on 1 Apr. for No. 134), the number of years and months during which contributions were paid in or widows' and orphans' pensions paid out, and also those that elapsed up to the date from which the pension is reckoned. From these data, the values in the adjacent columns have been determined. At the time of entry, for interest rates of \(3\tfrac{1}{2}\) per cent and 4 per cent, they are equivalent to the contributions paid annually at 10 Thl. and the widows' and orphans' pensions paid at 250 Thl. The amounts of 10 Thl. and 250 Thl. were probably chosen by \textsc{Gauss} because most of the table had already been calculated on this basis in 1845. For the widows still living on 1 Oct. 1851, namely those of members Nos. 86. 103. 104. 109. 112. 114. 122 and 134, the capital value of the pension payments still to be made after that date has been determined with the aid of \textsc{Brune}'s mortality tables, on the assumption that payment is made at the end of each year, but only through the end of the month of death. The mean capital values of contributions and pensions obtained directly from this compilation were used in the balance account of 1851, after multiplication by a correction factor to allow for pensions being paid half-yearly rather than annually. For orphans' pensions, it is assumed here that they continued until the end of the month in which the child completed its 20\textsuperscript{th} year of life.

Unlike the present printed version, the manuscript table has no headings for the individual columns.]
